% =============================================================================
%  Modelo LaTeX — Dissertação, Projeto, Relatório de Estágio ou Tese
%  Universidade Politécnica do Cávado e do Ave (UPCA), antigo IPCA
%
%  1. Escolha as opções e preencha os dados abaixo.
%  2. Escreva o texto nos ficheiros de Preambulo/, Capitulos/ e Anexos/.
%  3. Acrescente as referências em Bibliografia/referencias.bib.
%  Capa, folha de rosto, declaração, índices e listas são gerados automaticamente.
%  Instruções completas: README.md
% =============================================================================
\documentclass[
  dissertacao,  % dissertacao | projeto | estagio | tese
  pt,           % pt | en  (língua do trabalho; inglês requer anuência do orientador)
  upca,         % upca | ipca  (ipca = designação e capa do modelo oficial de 2018)
  provas,       % provas (entrega para defesa) | versaofinal (após a defesa)
  apa,          % apa (norma exigida) | ieee (só com autorização do orientador)
  twoside,      % twoside (frente e verso, como o modelo oficial) | oneside
  %integridade, % acrescenta uma declaração de integridade (art. 14.º do regulamento)
]{UPCAThesis}

% ------------------------------------------------------------------- Dados
\titulo{Título do Trabalho}            % em português (com a opção en, a capa usa o \tituloEN)
\subtitulo{Subtítulo do trabalho (se não existir, deixe vazio)}
\tituloEN{Title of the Work}            % em inglês (página do Abstract)
\subtituloEN{Subtitle of the work}
\autor{Nome Completo do Estudante}
\email{a00000@alunos.ipca.pt}            % opcional
\orientador{Prof. Doutor Nome do Orientador}   % \orientador[Orientadora]{...} no feminino
\coorientador{Prof. Doutor Nome do Coorientador}  % opcional
%\orientadorentidade{Eng. Nome -- Empresa}  % estágio: orientador na entidade de acolhimento
\curso{Engenharia Informática}           % "... Grau de Mestre em <curso>"
\data{}                                  % vazio = mês e ano atuais (p. ex. "outubro, 2026")
\palavraschave{palavra 1, palavra 2, palavra 3}  % até 5
\keywords{keyword 1, keyword 2, keyword 3}       % até 5
\reproducao{}           % integral | parcial | nenhuma — assinala a opção na declaração
%\assinatura{Imagens/assinatura.png}     % assinatura digitalizada (opcional)
%\declaracaoassinada{Anexos/declaracao.pdf} % substitui a declaração gerada pelo PDF impresso, assinado e digitalizado

\addbibresource{Bibliografia/referencias.bib}
% Siglas e abreviaturas: \newacronym{chave}{SIGLA}{Significado por extenso}
% No texto use \gls{chave}: na 1.ª vez aparece "Significado (SIGLA)", depois só "SIGLA".
% Só as siglas usadas entram na Lista de Abreviaturas e Siglas (por ordem alfabética).
\newacronym{apa}{APA}{American Psychological Association}
\newacronym{est}{EST}{Escola Superior de Tecnologia}
\newacronym{ipca}{IPCA}{Instituto Politécnico do Cávado e do Ave}
\newacronym{upca}{UPCA}{Universidade Politécnica do Cávado e do Ave}
                 % definição das siglas

% Pacotes adicionais
\usepackage{lipsum}\setlipsum{auto-lang=false} % só para o texto de exemplo; pode remover

\begin{document}

\paginasiniciais                         % capa, folha de rosto e declaração

\begin{apoios}
Nesta página faz-se alusão aos apoios financeiros recebidos (bolsas, projetos,
entidades financiadoras). Se não existirem, comente a linha
\verb|\begin{apoios}
Nesta página faz-se alusão aos apoios financeiros recebidos (bolsas, projetos,
entidades financiadoras). Se não existirem, comente a linha
\verb|\begin{apoios}
Nesta página faz-se alusão aos apoios financeiros recebidos (bolsas, projetos,
entidades financiadoras). Se não existirem, comente a linha
\verb|\input{Preambulo/Apoios}| no \texttt{MainThesis.tex}.
\end{apoios}
| no \texttt{MainThesis.tex}.
\end{apoios}
| no \texttt{MainThesis.tex}.
\end{apoios}
                 % opcional: comente (%) para remover
% O título, o subtítulo e as palavras-chave são inseridos automaticamente.
\begin{resumo}
O resumo deve ter entre 200 e 300 palavras e informar o leitor sobre a razão da
investigação, como foi feita, os principais resultados e as principais
conclusões. Esta secção é traduzida para inglês no Abstract.

\lipsum[1-2]
\end{resumo}
                 % trabalho em inglês: pode trocar a ordem
% Tradução do Resumo. O título em inglês e as keywords são inseridos automaticamente.
\begin{abstract}
The abstract should have between 200 and 300 words and is the English
translation of the Resumo.

\lipsum[3-4]
\end{abstract}
               % das duas linhas (Abstract primeiro)
\begin{agradecimentos}
Secção facultativa. Se não a utilizar, comente a linha
\verb|\begin{agradecimentos}
Secção facultativa. Se não a utilizar, comente a linha
\verb|\begin{agradecimentos}
Secção facultativa. Se não a utilizar, comente a linha
\verb|\input{Preambulo/Agradecimentos}| no \texttt{MainThesis.tex}.
\end{agradecimentos}
| no \texttt{MainThesis.tex}.
\end{agradecimentos}
| no \texttt{MainThesis.tex}.
\end{agradecimentos}
         % opcional
% Secção facultativa: o título não é impresso, mas aparece no índice.
\begin{dedicatoria}
Texto da dedicatória.
\end{dedicatoria}
            % opcional

\indices                                 % siglas, índice, figuras e tabelas

% Capítulos com \chapter* (como este) não são numerados, tal como no modelo oficial.
\chapter*{Introdução}

A introdução deve conter o enquadramento, a motivação para a escolha da
problemática, os objetivos do estudo, a metodologia a seguir e a estrutura do
trabalho.

\section{Estrutura do documento}

\lipsum[1]

\chapter{Exemplo de Capítulo}
\label{cap:exemplo}

Os títulos seguem automaticamente a formatação oficial. Qualquer secção com
subsecções deve ter pelo menos duas.

\section{Citações e siglas}

As referências seguem as normas da \gls{apa} \parencite{APA2020}. Pode citar
no texto, como \textcite{Deterding2011}, ou entre parênteses
\parencite{Resnick2009}. Na segunda utilização, a sigla aparece só na forma
curta (\gls{apa}). As notas de rodapé usam \verb|\footnote|.\footnote{Exemplo de nota.}

\begin{citacao}
Citações com mais de 40 palavras ficam em parágrafo próprio, em itálico e com
recuo de ambos os lados \parencite[p.~171]{APA2020}. \lipsum[2][1-3]
\end{citacao}

\section{Níveis de título}

\subsection{Título de nível 3}

\lipsum[3][1-4]

\subsubsection{Título de nível 4}

\lipsum[4][1-4]

\paragraph{Título de nível 5}

\lipsum[5][1-4]

\paragraph{Outro título de nível 5}

\lipsum[6][1-4]

\subsection{Outro título de nível 3}

Listas:
\begin{itemize}
  \item primeiro elemento;
  \item segundo elemento.
\end{itemize}

\chapter{Figuras, Tabelas e Código}

\section{Figuras}

A \autoref{fig:exemplo} mostra como inserir uma figura. Guarde as imagens na
pasta \texttt{Imagens/}.

\begin{figure}[htbp]
  \centering
  \includegraphics[width=0.5\linewidth]{example-image}
  \caption{Exemplo de figura (a legenda fica por baixo)}
  \label{fig:exemplo}
  \fonte{Elaboração própria.}
\end{figure}

\section{Tabelas}

Nas tabelas a legenda fica por cima (\autoref{tab:exemplo}).

\begin{table}[htbp]
  \centering
  \caption{Exemplo de tabela}
  \label{tab:exemplo}
  \begin{tabular}{lcc}
    \toprule
    Método & Precisão (\%) & Tempo (s) \\
    \midrule
    A & 91,2 & 1,4 \\
    B & 94,7 & 2,9 \\
    \bottomrule
  \end{tabular}
  \fonte{Adaptado de \textcite{Resnick2009}.}
\end{table}

\section{Equações e código}

\begin{equation}
  E = mc^2
  \label{eq:exemplo}
\end{equation}

\begin{lstlisting}[language=Python, caption={Exemplo de código}, label={lst:exemplo}]
def media(valores):
    # Calcula a média (comentários com acentuação são suportados)
    return sum(valores) / len(valores)
\end{lstlisting}

\chapter*{Conclusões}

Apresentam-se os aspetos centrais do trabalho, relacionam-se os objetivos
enunciados com os resultados obtidos e referem-se os contributos, as limitações
e o trabalho futuro.

\lipsum[7]


\referencias                             % referências bibliográficas

\anexos                                  % opcional (ou \apendices)
% Cada \chapter{...} depois de \anexos fica "ANEXO A – ...", "ANEXO B – ...".
\chapter{Descrição do Anexo}
\label{anx:exemplo}

Conteúdo do anexo.

%\anexopdf{Protocolo de estágio}{Anexos/protocolo.pdf}  % anexo que já existe em PDF

\end{document}
